%; whizzy chapter
% -initex iniptex -latex platex -format platex -bibtex jbibtex -fmt fmt
% 以上 whizzytex を使用する場合の設定。


%     Tokyo Debian Meeting resources
%     Copyright (C) 2006 Junichi Uekawa

%     This program is free software; you can redistribute it and/or modify
%     it under the terms of the GNU General Public License as published by
%     the Free Software Foundation; either version 2 of the License, or
%     (at your option) any later version.

%     This program is distributed in the hope that it will be useful,
%     but WITHOUT ANY WARRANTY; without even the implied warranty of
%     MERCHANTABILITY or FITNESS FOR A PARTICULAR PURPOSE.  See the
%     GNU General Public License for more details.

%     You should have received a copy of the GNU General Public License
%     along with this program; if not, write to the Free Software
%     Foundation, Inc., 51 Franklin St, Fifth Floor, Boston, MA  02110-1301 USA

%   Pdf作成手順
% dvipdfmx debianmeetingresume200605.dvi
%  preview (shell-command (concat "xpdf " (replace-regexp-in-string "tex$" "pdf"(buffer-file-name)) "&"))
% 画像ファイルを処理するためにはebbを利用してboundingboxを作成。
%(shell-command "cd image200605; ebb *.png")

%%ここからヘッダ開始。

\documentclass[mingoth,a4paper]{jsarticle}
\usepackage[dvipdfmx]{graphicx}
\usepackage{fancybox}
\usepackage{longtable}
\usepackage{ascmac}	% 囲み (screen,itembox)
\usepackage{fancyvrb}   % 囲み Verbatim のために必要
\usepackage[dvipdfmx]{hyperref}
\usepackage{url}

%http://www.naney.org/diki/dk/hyperref.html
%日本語EUC系環境の時
\AtBeginDvi{\special{pdf:tounicode EUC-UCS2}}
%シフトJIS系環境の時
%\AtBeginDvi{\special{pdf:tounicode 90ms-RKSJ-UCS2}}

%% spacing の設定をする。外枠を減らす。
\setlength\headheight{0mm}
\setlength\topmargin{-20mm}
\setlength\headsep{0mm}
\setlength\topskip{3mm}
\setlength\maxdepth{4pt}
\setlength\columnsep{6mm}
\setlength\textheight{252mm}
\setlength\topmargin{-5mm}
\setlength\textwidth{170mm}
\setlength\oddsidemargin{-5mm}
\setlength\evensidemargin{-5mm}

% commandline環境を定義。画面入出力についてはcommandline環境
% で表記する
\newenvironment{commandline}%
{\VerbatimEnvironment
  \begin{Sbox}\begin{minipage}{15cm}\begin{fontsize}{7.3}{7.3} \begin{BVerbatim}}%
{\end{BVerbatim}\end{fontsize}\end{minipage}\end{Sbox}
  \setlength{\fboxsep}{8pt}\fbox{\TheSbox}}


%%% start of santaku
\makeatletter
\newwrite\tf@jqz
\immediate\openout\tf@jqz\jobname.jqz\relax
\makeatother
\newcounter{santakucounter}
\newcommand{\santaku}[5]{%
\addtocounter{santakucounter}{1}

\addtocontents{jqz}{\arabic{santakucounter}. #5\\}
\nopagebreak 問題\arabic{santakucounter}. 
#1\\
\nopagebreak□ A #2\\
\nopagebreak□ B #3\\
\nopagebreak□ C #4
\pagebreak[1]
\hspace{1cm}
\\

}
%%% end of santaku

\newcommand{\emptyspace}{(\underline{\hspace{1cm}})}

\newcommand{\subsubsubsection}[1]{%
\vspace{1zw}{\bf #1}\\}

% sectionをセンタリングする
\makeatletter
  \renewcommand{\section}{\@startsection{section}{1}{\z@}%
    {\Cvs \@plus.5\Cdp \@minus.2\Cdp}% 前アキ
    {.5\Cvs \@plus.3\Cdp}% 後アキ
    {\normalfont\Large\headfont\raggedright\centering}} % style
\makeatother

% section の代わりの環境
\newcommand{\dancersection}[2]{%
\newpage
東京エリアDebian勉強会 2006
\hrule
\vspace{0.5mm}
\hrule
\hfill{}\includegraphics[width=3cm]{image200502/openlogo-nd.eps}\\
\vspace{-4cm}
\begin{center}
  \section{#1}
\end{center}
\hfill{}#2\hspace{3cm}\space\\
\hrule
\hrule
\vspace{1cm}
}

% for dancerj
\newcommand{\fgref}[1]{図\ref{#1}}
\newcommand{\tbref}[1]{表\ref{#1}}


\begin{document}

\begin{titlepage}

% 毎月変更する部分, 本文の末尾も修正することをわすれずに
\title{
 第16回 東京エリア Debian 勉強会\\事前資料}
\date{2006年5月14日}
\author{Debian勉強会会場係 上川 純一\thanks{Debian Project Official Developer}} 
\maketitle
\thispagestyle{empty}
\end{titlepage}

\newpage
\tableofcontents

\dancersection{Introduction To Debian 勉強会}{上川 純一}

今月のDebian勉強会へようこそ。
これからDebianのあやしい世界に入るという方も、すでにどっぷりとつかってい
るという方も、月に一回Debianについて語りませんか？

目的として下記の二つを考えています。

\begin{itemize}
 \item メールではよみとれない、もしくはよみとってられないような情報を情
       報共有する場をつくる
 \item まとまっていないDebianを利用する際の情報をまとめて、ある程度の塊と
       して出してみる
\end{itemize}

また、東京にはLinuxの勉強会はたくさんありますので、Debianに限定した勉強
会にします。Linuxの基本的な利用方法などが知りたい方は、他でがんばってくださ
い。
Debianの勉強会ということで究極的には参加者全員がDebian Packageを
がりがりと作りながらスーパーハッカーになれるような姿を妄想しています。

Debianをこれからどうするという能動的な展開への土台としての空間を提供し、
情報の共有をしたい、というのが目的です。
次回は違うこと言ってるかもしれませんが、御容赦を。

\subsection{講師紹介}

\begin{itemize}
 \item{上川 純一} 宴会の幹事です。
\end{itemize}

\subsection{事前課題紹介}

今回の事前課題は
「Debconf参加者に質問したいこと」
というタイトルで200-800文字程度の文章を書いてください。
というものでした。
その課題に対して下記の内容を提出いただきました。

\subsubsection{やまねさん}


* Debian はソフトウェアの品質を保つのにどのような努力をしていくのか？

リリースしたソフトを多数のユーザで利用してバグを報告→修正というサイクル
を行っているのが現在の方法だが、そもそもリリース時にバグを出さない方が
重要だ。無限にユーザの時間がある訳では無い。また、パッケージが増えていけ
ば、チェックされないことすらあるパッケージが増える。

パッケージングについては lintian check などがあるが、upstream のソース
が酷いようであれば、それを Debian でのパッケージにするまえに修正すべきだ。

テスト品質を明示的な一定基準に保つ方策を考えてほしい。
テストスイートを用意してコミット時にチェックなどする予定はないか。



* ソースへの標準的なアクセスを提供する予定は無いのか

開発中のパッケージのソースコードがバラバラに
パッケージとしてリリースされたソフトについては apt-get で取得できるが、


* Web を全面的に改定できないか

今の debian.org のサイトは必要な情報を見付けられない状態に陥っている。
改善の予定は無いのか


* ftpmaster の New Queue 処理

努力をしているのはよくわかるが、実際問題として受け入れまでが2週間程度
かかるのは遅すぎる（もし、これが遅いと思っていないのならその考え方が
一番の問題だが）。どのようにして改善するのか、解決策を明示し、実行を
宣言してほしい。

また、私が upload したパッケージは3週間をすぎてまだ放置されている。
問題があるなら直したいが、何が問題なのかすら ftpmaster からの連絡が
無い。せめてコメントを出すなど対処をしてはもらえないのだろうか。



\subsubsection{北原さん}


宿題：Debconf参加者に質問したいこと
回答：
この質問をされるのは、(1) 日本からDebconfに参加する人なのか、
(2) 会場にいる任意の（又は特定の）参加者に対してなのか、よく解らない
のですが、（(1)のような気がしますがあえて）(2)と仮定して聞きます。

\begin{itemize}
 \item  Q.1 Debconfが日本で開催されるのはいつ頃でしょうか？
 \item  Q.2 東京・京都以外で開催されるとすればどこがいいでしょうか？
 （２都市を外したのは、過去多くの国際会議が開催され面白くないから）
 \item  Q.3 Debconf前後に観光するとしたら、どこに行きたいですか？
\end{itemize}

Debian全然関係ないですね。


\dancersection{最近のDebian関連のミーティング報告}{上川 純一}

\subsection{東京エリアDebian勉強会15回目報告}
% (query-replace-regexp "<.*>" "")
	  東京エリアDebian勉強会報告。
	  4月の第15回Debian勉強会を実施しました。
	    今回は岩松さんがdebian-policyのソースの管理の部分について解説しました。
	    その後上川が日本語 TeX の Debian においての現状について調査報告しました。
	  
	  今回の参加人数は15人でした。
        
	
	  
	    Debian weekly news quizは小林さんが満点をとりました。おめでとうございます。
	  
	  
	    Debian TeX Policy について上川が解説しました。
	    最近 TeTeX と TeXLive が共存できるように変更が行われたり、
	    設定ファイルやディレクトリ構造の整理が行われているため、既存のパッケージが機能しなくなるという事象について解説しました。
	    また、日本語の処理できるTeXについてはplatex以外にもたくさんあることになっているが、
	    現状インストールしただけで利用できる状態になっているのは platex くらいだということを調査報告しました。
	  
	  
	    xpdfで日本語のしおりが出ない件についてすこし実演しました。説明しました。
	    現状、xpdfでは、しおりの部分がlatin-1以外の文字を削除している処理をしている、
	    日本語ではその部分はUTF-16でエンコードされており、それをEUC-JPなどに内部的に変換すればよいのだが、
	    表示の部分でMotifでの表示のさせかたがよくわからない。
	    gpdfでは、しおりの部分は表示されるのだが本文の文字が全く表示されない、という現状を説明しました。
	  
	  
	    岩松さんが、Debian policyについて紹介しました。
	    ソースパッケージについての解説をしました。
	    debian/以下のファイルがどうなっているのか、という細かい話しを解説しました。
	  
	    宴会は「一汁一菜　遇」にて開催。

\dancersection{次回}{}

6月17日です。
内容は本日決定予定です。

参加者募集はまた後程。

\newpage

\vspace*{15cm}
\hrule
\vspace{2mm}
\includegraphics[width=2cm]{image200502/openlogo-nd.eps}
\noindent \Large \bf Debian 勉強会資料\\ \\
\noindent \normalfont 2006年5月14日 \hspace{5mm}  初版第1刷発行\\
\noindent \normalfont 東京エリア Debian 勉強会 （編集・印刷・発行）\\
\hrule

\end{document}
